\documentclass[11pt, a4paper]{article}
\usepackage[a4paper, margin=2.5cm]{geometry}
\usepackage{fontspec}
\usepackage[portuguese, english, bidi=basic, provide=*]{babel}
\babelprovide[import, onchar=ids fonts]{portuguese}
\babelfont{rm}{Noto Sans}
\usepackage{enumitem}
\usepackage{xcolor}
\usepackage{listings}
\usepackage{tcolorbox}
\usepackage[hidelinks]{hyperref}

% Definicao de cores do documento
\definecolor{codebg}{RGB}{245, 247, 250}
\definecolor{codeframe}{RGB}{210, 215, 225}
\definecolor{primary}{RGB}{15, 76, 129}
\definecolor{accent}{RGB}{40, 116, 166}
\definecolor{warnbg}{RGB}{253, 237, 237}
\definecolor{warnframe}{RGB}{220, 53, 69}
\definecolor{infobg}{RGB}{240, 248, 255}
\definecolor{infoframe}{RGB}{40, 116, 166}

% Linguagem YAML personalizada para listings
\lstdefinelanguage{yaml}{
  keywords={true, false, null, y, n},
  keywordstyle=\color{primary}\bfseries,
  basicstyle=\ttfamily\small,
  sensitive=false,
  comment=[l]{\#},
  commentstyle=\color{gray}\itshape,
  stringstyle=\color{accent},
  morestring=[b]',
  morestring=[b]"
}

% Configuracao basica do listings
\lstset{
  backgroundcolor=\color{codebg},
  frame=single,
  rulecolor=\color{codeframe},
  basicstyle=\ttfamily\small,
  breaklines=true,
  columns=fullflexible,
  keepspaces=true,
  showstringspaces=false,
  commentstyle=\color{gray}
}

\setlength{\parindent}{0pt}
\setlength{\parskip}{7pt}

\begin{document}

\begin{center}
    {\LARGE \textbf{\color{primary} Laboratório de Kubernetes}} \\[4pt]
    {\large \textbf{\color{accent} Da Máquina Local à Nuvem Híbrida}} \\[8pt]
    \rule{\textwidth}{1pt}
\end{center}

\textbf{Objetivo:} Sair do zero, criar um cluster Kubernetes local, testar a auto-recuperação e, por fim, simular uma arquitetura de nuvem híbrida.

\begin{tcolorbox}[colback=infobg, colframe=infoframe, arc=3pt, boxrule=0.8pt]
\textbf{INSTRUÇÕES DE ENTREGA DA ATIVIDADE:}
\begin{enumerate}[leftmargin=*, topsep=3pt, itemsep=2pt]
    \item Execute todos os comandos detalhados neste roteiro prático.
    \item Capture telas (\textit{prints}) que comprovem a execução e os resultados dos comandos no seu terminal.
    \item Elabore e entregue um \textbf{relatório individual} explicando o que foi feito em cada etapa, ilustrado com os seus respectivos \textit{prints}.
\end{enumerate}
\end{tcolorbox}

\begin{tcolorbox}[colback=warnbg, colframe=warnframe, arc=3pt, boxrule=0.8pt]
\textbf{ATENÇÃO:} Em todos os comandos deste guia, usaremos o prefixo \texttt{sudo k3s kubectl}.
\end{tcolorbox}

\section*{Fase 1: O Básico (Comandos e Deploy)}

Vamos instalar o cluster e criar a nossa primeira aplicação.

\subsection*{1. Instalar o Kubernetes (k3s)}
No seu terminal, execute o comando de instalação do k3s:

\begin{lstlisting}[language=bash]
curl -sfL https://get.k3s.io | sh -
\end{lstlisting}

Verifique se o nó (sua máquina) está pronto:

\begin{lstlisting}[language=bash]
sudo k3s kubectl get nodes
\end{lstlisting}

\subsection*{2. Criar um Deployment (Modo Declarativo)}
Vamos criar um gerenciador (\texttt{Deployment}) que garante 3 réplicas do Nginx:

\begin{lstlisting}[language=bash]
sudo k3s kubectl create deployment meu-app --image=nginx:latest --replicas=3
\end{lstlisting}

\subsection*{3. Expor e Testar}
Exponha a aplicação e veja o serviço criado:

\begin{lstlisting}[language=bash]
sudo k3s kubectl expose deployment meu-app --type=NodePort --port=80
sudo k3s kubectl get services
\end{lstlisting}

\textit{Anote a porta que apareceu no \texttt{PORT(S)} (ex: \texttt{80:31218/TCP}) e acesse no navegador: \texttt{http://localhost:31218}.}

\section*{Fase 2: Testando Conceito Local (Auto-cura)}

Vamos provar que o Deployment mantém o estado desejado.

1. \textbf{Simule uma falha:} \\
Pegue o nome de um dos pods listados abaixo (será algo como \texttt{meu-app-xxxxxxxx-xxxxx}):

\begin{lstlisting}[language=bash]
sudo k3s kubectl get pods
\end{lstlisting}

Delete o pod (substitua pelo nome que você copiou):

\begin{lstlisting}[language=bash]
sudo k3s kubectl delete pod meu-app-<hash>
\end{lstlisting}

\textbf{2. Observando a mágica:} \\
Rode o comando com \texttt{-w} (\textit{watch}) para ver em tempo real:

\begin{lstlisting}[language=bash]
sudo k3s kubectl get pods -w
\end{lstlisting}

\begin{tcolorbox}[colback=infobg, colframe=infoframe, arc=3pt, boxrule=0.8pt]
\textbf{OBSERVE:} O pod que você deletou entrou em \texttt{Terminating} e um \textbf{novo pod foi criado automaticamente} em segundos. O Kubernetes percebeu que você queria 3 e recriou o que faltava. \textit{(Aperte \texttt{Ctrl+C} para sair do watch).}
\end{tcolorbox}

\subsection*{3. LIMPEZA (A Transição)}
Agora que já entendemos o básico, vamos apagar o \texttt{meu-app} para que a nossa próxima aplicação (a distribuída na nuvem) seja a única a rodar o Nginx no cluster. Isso evita confusão na hora de olhar os pods.

\begin{lstlisting}[language=bash]
sudo k3s kubectl delete deployment meu-app
sudo k3s kubectl delete service meu-app
\end{lstlisting}
\textit{(Verifique com \texttt{sudo k3s kubectl get pods} se não sobrou nada).}

\section*{Fase 3: A Nuvem Híbrida (Simulando a AWS)}

Agora, vamos colocar uma réplica da sua aplicação no ``datacenter do professor''.

\subsection*{1. Pegar seus dados de conexão}
Pegue o IP da sua máquina e o Token do seu cluster:

\begin{lstlisting}[language=bash]
hostname -I
sudo cat /var/lib/rancher/k3s/server/node-token
\end{lstlisting}

\subsection*{2. Contratar a Nuvem}
Abra o navegador e acesse o portal do professor: \texttt{http://<IP\_DO\_PROFESSOR>:1961}. \\
Preencha o seu nome, o IP da sua máquina e o Token. Clique em \textbf{Conectar}.

\subsection*{3. Verificar a conexão}
Rode:

\begin{lstlisting}[language=bash]
sudo k3s kubectl get nodes
\end{lstlisting}

\begin{tcolorbox}[colback=warnbg, colframe=warnframe, arc=3pt, boxrule=0.8pt]
\textbf{ATENÇÃO:} Agora você verá \textbf{dois} nós: o seu computador (\textit{Control Plane}) e um nó com um nome tipo \texttt{2945794d5ccd} ou \texttt{worker-seu-nome}. \textbf{Anote o nome do nó local e o nome do nó da nuvem}, você precisará deles!
\end{tcolorbox}

\subsection*{4. Deploy Distribuído (O Esqueleto YAML)}
Vamos criar uma aplicação com a imagem \texttt{nginxdemos/hello}, que mostra o nome do pod que respondeu à requisição.

Crie um arquivo chamado \texttt{app-nuvem.yaml} e \textbf{preencha os campos em branco}:

\begin{lstlisting}[language=yaml]
apiVersion: apps/v1
kind: Deployment
metadata:
  name: app-distribuida
spec:
  # PREENCHA AQUI: Quantas replicas queremos? Lembre-se, queremos que existam replicas na nuvem e local.
  replicas: 
  selector:
    matchLabels:
      app: distribuida
  template:
    metadata:
      labels:
        app: distribuida
    spec:
      # =========================================================================
      # EXERCÍCIO DE ANTI-AFFINITY (O Pulmão na Nuvem)
      # =========================================================================
      # Objetivo: Impedir que o Kubernetes coloque dois Pods desta aplicação
      # na mesma máquina física.
      # =========================================================================
      affinity:
        podAntiAffinity:
          # Usamos "required" (Obrigatório) para forçar a separação. Se a máquina
          # local estiver bloqueada (cordon), o pod ficará "Pending" em vez de
          # desrespeitar a regra.
          requiredDuringSchedulingIgnoredDuringExecution:
- labelSelector:
                matchExpressions:
                  # -----------------------------------------------------------
                  # PREENCHA AQUI: Qual é a chave (key) da label do nosso Pod?
                  # -----------------------------------------------------------
                  # COMO DESCOBRIR:
                  # Rode o comando abaixo no terminal para ver o manifesto
                  # do Deployment que você acabou de criar:
                  #
                  #   sudo k3s kubectl get deployment app-distribuida -o jsonpath='{.spec.selector.matchLabels}'
                  #
                  # O resultado será algo como {"app":"distribuida"}.
                  # A "key" é o que está antes dos dois pontos ("app").
- key: [PREENCHA COM A CHAVE DA LABEL]

                    operator: In

                    # -----------------------------------------------------------
                    # PREENCHA AQUI: Qual é o valor (value) dessa label?
                    # -----------------------------------------------------------
                    # COMO DESCOBRIR:
                    # No mesmo resultado do comando anterior, o "value" é
                    # o que está depois dos dois pontos ("distribuida").
                    values: [[PREENCHA COM O VALOR DA LABEL]]

              # ---------------------------------------------------------------
              # PREENCHA AQUI: Qual é a chave que representa uma "máquina física"?
              # ---------------------------------------------------------------
              # COMO DESCOBRIR:
              # 1. Pesquise na documentação do Kubernetes sobre "topologyKey".
              # 2. Ou rode este comando para ler a explicação do próprio K8s:
              #
              #   sudo k3s kubectl explain pod.spec.affinity.podAntiAffinity.requiredDuringSchedulingIgnoredDuringExecution.topologyKey
              #
              # O texto dirá que esta chave define o "domínio de topologia".
              # A label padrão para identificar uma máquina física (Nó) é:
              #   kubernetes.io/hostname
              topologyKey: [PREENCHA COM A CHAVE DE TOPOLOGIA]
      containers:
- name: web
          # PREENCHA AQUI: Use a imagem que mostra o nome do Pod.
          # DICA: Pesquise por "nginxdemos/hello" no Docker Hub (hub.docker.com). 
          # A versão que queremos é a "plain-text", que retorna texto simples.
          image: 
          ports:
- containerPort: 80
---
apiVersion: v1
kind: Service
metadata:
  name: app-distribuida
spec:
  type: NodePort
  selector:
    app: distribuida
  ports:
- port: 80
      targetPort: 80
      # PREENCHA AQUI: Fixe a porta externa em 30080 para todos acessarem igual
      nodePort: 
\end{lstlisting}

Aplique o arquivo e veja onde os pods caíram:

\begin{lstlisting}[language=bash]
sudo k3s kubectl apply -f app-nuvem.yaml
sudo k3s kubectl get pods -o wide
\end{lstlisting}

\begin{tcolorbox}[colback=infobg, colframe=infoframe, arc=3pt, boxrule=0.8pt]
\textbf{OBSERVE:} Na coluna \texttt{NODE}, um pod deve estar no \textbf{seu nó local} e o outro no \textbf{nó da nuvem}. Anote o nome exato do pod que caiu no seu nó local!
\end{tcolorbox}

\subsection*{5. O Teste de Fogo (Derrubando o Local)}
Vamos provar que a nuvem segura a bronca. Vamos simular que o seu datacenter local caiu.

1. \textbf{Deixe o terminal observando os pods:}
\begin{lstlisting}[language=bash]
sudo k3s kubectl get pods -o wide -w
\end{lstlisting}

\textbf{2. Em outro terminal, isole o seu nó local:} \textit{(Use o nome do nó local anotado)}
\begin{lstlisting}[language=bash]
sudo k3s kubectl cordon <NOME_DO_SEU_NO_LOCAL>
\end{lstlisting}

\textbf{3. Simule a queda do pod local:} \textit{(Use o nome do pod local anotado)}
\begin{lstlisting}[language=bash]
sudo k3s kubectl delete pod <NOME_DO_POD_LOCAL>
\end{lstlisting}

\begin{tcolorbox}[colback=infobg, colframe=infoframe, arc=3pt, boxrule=0.8pt]
\textbf{O QUE ACONTECEU?} Olhe para o terminal do \textit{watch}. O pod que você deletou vai tentar renascer, mas vai ficar \textbf{\texttt{Pending}} (Pendente). O Kubernetes está protegendo a sua regra de \textit{Anti-affinity}: ele prefere deixar o pod pendente a colocar dois pods na nuvem.
\end{tcolorbox}

\subsection*{6. A Prova Final (O Terminal não mente)}
Deixe o \texttt{curl} rodando em loop para provar que o site não caiu:

\begin{lstlisting}[language=bash]
while true; do curl -s http://localhost:30080 | grep "Server name"; sleep 1; done
\end{lstlisting}

\begin{tcolorbox}[colback=infobg, colframe=infoframe, arc=3pt, boxrule=0.8pt]
\textbf{OBSERVE:} O terminal vai mostrar o nome do pod da nuvem respondendo. O site \textbf{continua no ar}, servido 100\% pelo datacenter do professor! \textit{(Aperte \texttt{Ctrl+C} para parar).}
\end{tcolorbox}

\subsection*{7. Restaurando o Ambiente}
Para voltar tudo ao normal e o pod \texttt{Pending} voltar para a sua máquina:

\begin{lstlisting}[language=bash]
sudo k3s kubectl uncordon <NOME_DO_SEU_NO_LOCAL>
\end{lstlisting}

\section*{Limpeza Final}

Para não sobrecarregar o PC do professor, apague a aplicação da nuvem e destrua o seu cluster local:

\begin{lstlisting}[language=bash]
sudo k3s kubectl delete deployment app-distribuida
sudo k3s kubectl delete service app-distribuida
/usr/local/bin/k3s-uninstall.sh
\end{lstlisting}

\end{document}